\section{Comparison with ordinary combinations}\label{sec:comparison}

\Cref{thm:square} supplies the square estimate needed by
\cref{prop:comparison} at \(p=2.258\).
We retain the general statement for any available square estimate
\(w(1)\le p\) with \(2<p\le2.258\), together with the specified
baseline and the hypotheses on further inputs below.
We also give a comparison through aspect ratio \(0.93\) using only
the finite-family witnesses of \cref{sec:certificates} on the
new-algorithm side.

We compare upper-bound estimates, not unknown exact values of \(W\).
A \emph{cost tuple} \((a,b,c;e)\) records the certified estimate
\(W(a,b,c)\le e\).
The operations considered below act on such recorded estimates.

\begin{definition}\label{def:baseline}
For \(p>2\), the table baseline \(\mathcal B_p\) is generated by:
\begin{enumerate}
 \item the numerical entries of Table~1 of \cite{vxxz,advxxz}, with the
 exponent-two endpoint entered as \((1,0.321334,1;2)\);
 \item the square input \((1,1,1;p)\);
 \item the naive matrix-multiplication estimates.
\end{enumerate}
We close these inputs under permutations, scaling, tensor products,
padding, ordinary blocking and convex interpolation.
Ordinary subdivisions into products covering an index box may also
be used.
More precise input values may be substituted if their separator tests
below are verified.
\end{definition}

We use Table~1 of arXiv:2307.07970v2 and arXiv:2404.16349v3.
The latter's version~2 has the same listed rectangular entries;
the updated square entry does not affect the argument.
This definition makes the numerical baseline reproducible.
It does not assert that the tables exhaust every possible parameter
choice in the source algorithms.

\subsection{A separator preserved by the baseline operations}

Let
\begin{equation}\label{eq:plane-parameters}
 q_*=0.45,\qquad
 v_p=\frac{p-2}{1-q_*},\qquad
 u_p=1-\frac{q_*v_p}{2},\qquad 2<p\le2.258.
\end{equation}
Then \(0<v_p\le u_p\le1\).
For a triple with smallest coordinate \(c\), put
\begin{equation}\label{eq:plane}
 P_p(a,b,c)=u_p(a+b)+v_pc.
\end{equation}
Equivalently, without ordering the coordinates,
\[
 P_p(a,b,c)
 =u_p(a+b+c)+(v_p-u_p)\min(a,b,c),
\]
the maximum of the three linear forms that assign \(v_p\) to one
coordinate and \(u_p\) to the other two.

\begin{lemma}\label{lem:plane-preservation}
The condition \(e\ge P_p(a,b,c)\) is preserved by all the operations
in \cref{def:baseline}.  The square input satisfies it with equality,
and the naive estimates satisfy it.
\end{lemma}

\begin{proof}
Each of the three linear-form tests adds under tensor products,
positive scaling and convex interpolation; this proves the condition
for the maximum as well.
It is invariant under permutations and monotone in the dimensions.
Padding therefore preserves it.
Increasing one dimension's exponent by \(d\) changes \(P_p\) by at
most \(d\), because all slopes lie in \([0,1]\); ordinary blocking
adds \(d\) to the cost.

The corresponding weighted-box inequality also holds.
For a box of side lengths \(L_i\), normalize the side lengths of a
contained subbox to ratios \(r_i\in[0,1]\).
For any of the weight triples \((u_p,u_p,v_p)\),
\(\prod_i r_i^{w_i}\ge\prod_i r_i\).
If subboxes cover the original box, their volume ratios sum to at
least one.  Hence the weighted size of the whole box is at most
the sum of the weighted sizes of the covering boxes.
Intersecting with the original box first handles boxes extending
outside it.
Thus summing declared costs for ordinary subdivided products cannot
cross the separator either.
Precisely, at scale \(n\) the tuple \((a,b,c;e)\) declares a charge
\(n^e\) for a box with side lengths \(n^a,n^b,n^c\).
Apply the covering inequality to these charges before taking exponent
limits; it remains valid for subdivisions depending on \(n\).

Finally
\(2u_p+v_p=2+(1-q_*)v_p=p\), and all weights are at most one.
These give the square and naive assertions.
\end{proof}

For any additional old input it suffices to check the plane at
\(p=2.258\) together with the usual size bound.
Indeed, after ordering,
\[
 P_p(a,b,c)=a+b+v_p\bigl(c-q_*(a+b)/2\bigr).
\]
If the parenthesis is nonnegative, this increases with \(p\).
Otherwise it is at most the size bound \(a+b\).

For shapes \((1,\ell,1)\), the test at \(p=2.258\) is
\begin{equation}\label{eq:shape-plane}
 P_{2.258}(1,\ell,1)=
 \begin{cases}
  2+\dfrac{0.258}{0.55}(\ell-0.45),&0\le\ell\le1,\\[2mm]
  u_{2.258}(1+\ell)+v_{2.258},&\ell\ge1.
 \end{cases}
\end{equation}
If \(\ell\le0.45\), or \(\ell\ge2/0.45-1\), it already follows
from the size bound.
Between those regimes the following downward roundings of the
published entry values suffice.  They are bounds on the \emph{reported
costs}, not lower bounds on matrix-multiplication exponents.
\[
\begin{array}{c|rrrrrrrr}
\ell&.50&.527\text{--}.528&.55&.60&.65,.70&
 .75,.80&.85,.90,.95&1\\ \hline
\text{entry at least}&2.04&2.05&2.06&2.09&2.12&2.18&2.24&2.36
\end{array}
\]
\[
\begin{array}{c|rrrrr}
\ell&1.10,1.20&1.50&2&2.50&3\\ \hline
\text{entry at least}&2.44&2.71&3.16&3.60&4.05.
\end{array}
\]
For example, the newer entries at \(\ell=.50,.60,.70\) are
\(2.042776,2.092351,2.152770\), respectively \cite{advxxz}.
All displayed roundings satisfy \eqref{eq:shape-plane} by substitution.
The other entries either lie in the automatic size-bound region or
are covered by these grouped tests.
It follows from \cref{lem:plane-preservation} that every estimate
generated by \(\mathcal B_p\) remains on or above \(P_p\).

\subsection{The strict improvement interval}

\begin{proposition}\label{prop:comparison}
Assume \(2<p\le2.258\) and \(w(1)\le p\).
Against \(\mathcal B_p\), the new bounds are strictly better for
\(0.321334<k<1\).
More generally they improve any separator-satisfying baseline
throughout the part of \(0<k<1\) on which that baseline exceeds two.
\end{proposition}

\begin{proof}
For \(k\le0.465\), \eqref{eq:alpha-rational} gives \(w(k)=2\).
For \(0.465<k<1\), convexity and the available square point give
\[
 w(k)\le 2+\frac{p-2}{1-0.465}(k-0.465).
\]
The difference between the separator at \((1,k,1)\) and this line is
\begin{equation}\label{eq:positive-gap}
 \frac{(p-2)(1-k)(0.465-0.45)}
 {(1-0.45)(1-0.465)}>0.
\end{equation}
This proves the general assertion.

It remains to identify the exponent-two range of the specified table
baseline, including limits of ordinary combinations.
Put \(q_0=0.321334\).
For each of the finitely many table inputs, order its dimensions so
that \(c\) is smallest and write
\[
 d=c-q_0(a+b)/2,\qquad E=e-a-b.
\]
Every entry has \(E\ge0\), and every entry with \(d>0\) has
\(E>0\): the only table entries attaining the size bound have
\(2c/(a+b)\le q_0\).
Choose \(\eta>0\) such that
\[
 \eta\le\frac12,\qquad
 \eta\le\frac{p-2}{1-q_0},\qquad
 \eta\le\frac{E}{d}\quad\hbox{for every table entry with }d>0.
\]
Such a choice exists because the list is finite and all the relevant
ratios are positive.
Set \(u_0=1-q_0\eta/2\).
Then \(0<\eta\le u_0\le1\), and every table input satisfies
\[
 e\ge u_0(a+b)+\eta c=a+b+\eta d.
\]
For \(d\le0\) this follows from the size bound; for \(d>0\) it
follows from the choice of \(\eta\).
The square input satisfies the same test since
\(2u_0+\eta=2+(1-q_0)\eta\le p\); naive inputs also satisfy it.
The proof of \cref{lem:plane-preservation} applies verbatim to the
weights \((u_0,u_0,\eta)\), including its weighted-box argument for
subdivisions.
Every baseline estimate at \((1,k,1)\), and every limit of such
estimates, is therefore at least
\[
 2+\eta(k-q_0)>2\qquad(k>q_0).
\]
The recorded endpoint and padding give cost two for \(k\le q_0\).
Thus the baseline has flat endpoint exactly \(q_0\), and the general
assertion gives the claimed open interval.
\end{proof}

The endpoint at one is excluded because the comparison uses the same
square point on both sides.
By \cref{thm:square}, \(w(1)<2.258\), so
\cref{prop:comparison} applies at \(p=2.258\) and proves strict
improvement over \(\mathcal B_{2.258}\) throughout
\(0.321334<k<1\).
If a more precise old exponent-two endpoint is supplied, the same
statement uses that endpoint in place of the rounded table entry.
No strict improvement is asserted where the baseline already gives two.

There is also a comparison using only the constructions in
\(\calF_7\).  Interpolate their three proved points
\[
 (0.465,2),\qquad(0.709,2.092),\qquad(1,2.267).
\]
Against the numerical \(p=2.258\) separator, the comparisons at
\(k=0.465,0.709,0.93\) are
\[
\begin{array}{c|ccc}
k&0.465&0.709&0.93\\ \hline
\text{new line}&2&2.092&2.224903780\ldots\\
\text{separator}&2.007036363\ldots&2.121494545\ldots&
 2.225163636\ldots .
\end{array}
\]
Both differences are affine on the intervening segments.
The smallest endpoint margin is greater than \(0.0002598\);
these inequalities also follow directly from rational arithmetic.
Thus this finite family improves that numerical baseline through
\(k=0.93\), in addition to the flat-range improvement.

\subsection{A sufficient test for further old estimates}

The table comparison can be extended without evaluating a large
nonconvex parameter program, provided the extra input estimates
satisfy an explicit condition.
Let \(f\) be the convex upper-bound envelope from such old estimates
\emph{before} adding the new square point.
Suppose
\begin{equation}\label{eq:old-envelope-assumption}
 f(.50)\ge2.035,\qquad f(.60)\ge2.087,
\end{equation}
and include the known estimates
\[
 f(.40)\le2.010,\quad f(.50)\le2.043,\quad
 f(.60)\le2.093,\quad f(.70)\le2.154.
\]
Condition \eqref{eq:old-envelope-assumption} concerns the specified
envelope of estimates.  The published record values alone do not
prove it for every possible untabulated parameter choice.

\begin{proposition}\label{prop:extra-inputs}
Under \eqref{eq:old-envelope-assumption}, every such additional input
satisfies the separator test.
\end{proposition}

\begin{proof}
If an input \((a,b,c;e)\), with \(c\) smallest, violated the test,
average its two large axes and normalize by \((a+b)/2\).
It would give a violating balanced input at
\(x=2c/(a+b)\in[0,1]\).
The case \(x\le.45\) is excluded by the size bound.
For the rest, convexity extends a chord beyond its endpoints as
a lower bound.
Using a lower bound at the nearer endpoint and an upper bound at
the farther one gives
\[
\begin{array}{c|c|cc}
x\text{ interval}&\text{lower bound for }f(x)&
 \multicolumn{2}{c}{\text{gap above }2+\frac{.258}{.55}(x-.45)
                   \text{ at endpoints}}\\ \hline
{[.45,.50]}&2.035+.58(x-.50)&3/500&127/11000\\
{[.50,.55]}&2.035+.25(x-.50)&127/11000&13/22000\\
{[.55,.60]}&2.087+.67(x-.60)&29/4400&183/11000\\
{[.60,1]}&2.087+.44(x-.60)&183/11000&1/200 .
\end{array}
\]
The gap is affine and positive at both endpoints in every row,
contradicting the hypothetical violation.
\end{proof}

This is a sufficient inclusion rule for further baseline inputs, not
a claim about all future analyses of the source tensors.
For an enlarged baseline, \cref{prop:comparison} gives strict
improvement wherever that baseline exceeds two in \(0<k<1\).
The comparison with \(\mathcal B_{2.258}\) uses the proved
dual-exponent witness together with \cref{thm:square};
the extension in \cref{prop:extra-inputs} still requires
\eqref{eq:old-envelope-assumption}.
The finite-family numerical bounds were established independently
in \cref{sec:certificates}.

\paragraph{Scope of the result.}
The family yields asymptotic arithmetic exponents with arbitrarily
small positive slack.  It does not claim a practical crossover size
or global optimality of the chosen seven-completion family.
The finite-family entropy argument and parameter certificates are
independent of the square construction in \cref{thm:square}.
In the rectangular support bound and this comparison, an available
square estimate enters through its explicitly charged exponent point
\((2/p,1/p)\).

\section{Finite-exception characteristic transfer}\label{sec:characteristic-transfer}

For a field \(F\) and \(0\le k\le1\), let \(w_F(k)\) denote the arithmetic exponent for
the dimensions \((n,\lceil n^k\rceil,n)\), defined as in the
introduction but counting operations over \(F\) and allowing constants
from \(F\).  Put \(\omega_F=w_F(1)\) and
\(\alpha_F=\sup\{k\in[0,1]:w_F(k)=2\}\).  The constants in operation-count
bounds may depend on the field and the fixed schemes below; coefficient
bit sizes and numerical stability are not part of this model.

\begin{corollary}[Finite-exception field bounds]\label{cor:characteristic-transfer}
There is a finite set \(\mathcal S\) of positive primes such that, for
every field \(F\) with \(\operatorname{char}F=0\) or
\(\operatorname{char}F=p\notin\mathcal S\),
\[
 \omega_F<\frac{1129}{500}=2.258,\qquad
 w_F\!\left(\frac{709}{1000}\right)<\frac{523}{250}=2.092.
\]
The same set \(\mathcal S\) works for both inequalities and depends on
two fixed exact rank schemes selected in the proof.
\end{corollary}

\begin{proof}
First select the schemes over \(\C\).  The strict square and rectangular
bounds in Theorem~\ref{thm:main}, together with the rank limit in
Lemma~\ref{lem:rank-exponent}, give fixed exact rank schemes for
\(\langle m,m,m\rangle\) and \(\langle a,b,a\rangle\), with respective
rank budgets \(r_{\mathrm{sq}}\) and \(r_{\mathrm{rec}}\), such that
\begin{equation}\label{eq:fixed-transfer-witnesses}
 r_{\mathrm{sq}}^{500}<m^{1129},\qquad
 r_{\mathrm{rec}}^{250}<a^{523},\qquad
 r_{\mathrm{rec}}^{709}<b^{2092}.
\end{equation}
Indeed, choose finite values \(s_{\mathrm{sq}},s_{\mathrm{rec}}>0\)
at which the rank ratios in that lemma for \((1,1,1)\) and
\((1,709/1000,1)\) are respectively below \(1129/500\) and \(523/250\).
Take \(m=\lceil e^{s_{\mathrm{sq}}}\rceil\),
\(a=\lceil e^{s_{\mathrm{rec}}}\rceil\), and
\(b=\lceil e^{709s_{\mathrm{rec}}/1000}\rceil\).  Their lower bounds
by the unrounded quantities give \eqref{eq:fixed-transfer-witnesses}.
Thus the two strict margins concern only fixed positive integers;
no limiting degeneration is being specialized.

Write each rank-one term as a product of three linear forms and regard
all their coefficients in both schemes as indeterminates.  For
\(j\in\{\mathrm{sq},\mathrm{rec}\}\), let \(M_j\) be the corresponding
matrix tensor.  Coefficient matching consists of
finitely many equations
\[
 \sum_{\nu=1}^{r_j}
 \xi_{j,\nu,x}\eta_{j,\nu,y}\zeta_{j,\nu,z}
 =(M_j)_{xyz}\in\{0,1\}.
\]
The ideal \(I\) generated by these equations in a polynomial ring
\(\mathbb Q[\mathbf u]\) is proper because the selected complex
schemes give a joint solution.  Choose a maximal ideal containing
\(I\).  Its residue field \(K\) is a finite extension of \(\mathbb Q\)
by the Hilbert Nullstellensatz
\cite[Theorem~10.34.1, Tag~00FV]{stacks-nullstellensatz}.
The images of the indeterminates give both exact schemes over this
one number field.  Individual summands are allowed to vanish: only
the upper rank budgets matter, so no nonvanishing constraint is needed.

Fix a \(\mathbb Q\)-basis \(e_1=1,e_2,\ldots,e_d\) of \(K\).
Choose an integer \(D\ge1\) clearing the rational coordinates of all
selected coefficients in this basis and the structure constants of
all products \(e_i e_j\).  With \(R=\mathbb Z[1/D]\),
\[
 \mathcal A=\bigoplus_{i=1}^d R e_i\subset K
\]
is then a unital \(R\)-algebra, free of rank \(d\) as an \(R\)-module,
containing the coefficients of both identities.  Let \(\mathcal S\)
be the finite set of prime divisors of \(D\).

For any field \(F\) in the statement, the canonical map \(R\to F\)
is defined.  The algebra \(\mathcal A\otimes_R F\) is a nonzero
unital \(F\)-algebra of dimension \(d\).  A quotient by a maximal
ideal is therefore a finite field extension \(E/F\) of degree at
most \(d\).  The map \(\mathcal A\to E\) specializes both exact
tensor identities with their dimensions and rank budgets unchanged.
In positive characteristic this is reduction modulo that
characteristic followed by a finite extension; when \(F=\mathbb F_p\),
the quotient \(E\) is a finite residue field.

It remains to measure operations over \(F\), not \(E\).  The recursive
blocking and padding upper bound in the proof of
Lemma~\ref{lem:rank-exponent} uses only an exact rank scheme and is valid
over \(E\).  Set
\[
 \sigma=\min\left\{\log a,\frac{1000}{709}\log b\right\}>0.
\]
The integer inequalities \eqref{eq:fixed-transfer-witnesses} give
\[
 \frac{\log r_{\mathrm{sq}}}{\log m}<\frac{1129}{500},\qquad
 \frac{\log r_{\mathrm{rec}}}{\sigma}<\frac{523}{250}.
\]
For the rectangular scheme, depth \(q=\lceil\log n/\sigma\rceil\)
provides outer dimension at least \(n\) and inner dimension at least
\(n^{709/1000}\), at cost \(O((q+1)r_{\mathrm{rec}}^q)\) operations
over \(E\); the square scheme has the analogous bound with
\(q=\lceil\log n/\log m\rceil\).
Run each whole recursive computation on \(F\)-valued inputs inside
\(E\), and express every addition and multiplication in a fixed
\(F\)-basis of \(E\) containing \(1\).  The multiplication table
simulates each such operation with a fixed number of \(F\)-operations,
independent of \(n\); projection onto the coordinate of \(1\) recovers
the \(F\)-valued output.  This is a constant factor on the whole
operation count, not a new factor in the rank budget at every tensor
level.  The displayed strict exponent bounds therefore hold over \(F\).
\end{proof}

The argument neither computes \(\mathcal S\) nor covers its primes.
It applies only to the two fixed schemes chosen above.

\begin{corollary}[Characteristic-zero dual bound]\label{cor:characteristic-zero-dual}
For every field \(F\) of characteristic zero, \(\alpha_F>0.465\).
\end{corollary}

\begin{proof}
By Theorem~\ref{thm:main} and the definition of \(\alpha\), fix
\(k\in(0.465,1]\) with \(w(k)=2\).  For each \(\varepsilon>0\),
Lemma~\ref{lem:rank-exponent} gives a finite \(t>0\) and an exact
rank-\(r\) scheme for \(\langle a,b,a\rangle\), where
\[
 a=\lceil e^t\rceil,\qquad b=\lceil e^{kt}\rceil,\qquad
 \frac{\log r}{t}<2+\frac{\varepsilon}{2}.
\]
The coefficient-equation descent in the proof of
Corollary~\ref{cor:characteristic-transfer} realizes this finite scheme
over a number field \(K_{\varepsilon}\).  A maximal-ideal quotient of
\(K_{\varepsilon}\otimes_{\mathbb Q}F\) is a finite field extension
\(E/F\) over which the same exact rank budget holds.

Recursive blocking and padding at depth \(h=\lceil\log n/t\rceil\)
cover dimensions \((n,\lceil n^k\rceil,n)\).  As in
Lemma~\ref{lem:rank-exponent}, its operation count over \(E\) is
\[
 O((h+1)r^h)=O\!\left(\log n\,n^{\log r/t}\right)
             =O(n^{2+\varepsilon}).
\]
For this fixed accuracy, simulate the whole circuit in a fixed
\(F\)-basis of \(E\) containing \(1\), as in
Corollary~\ref{cor:characteristic-transfer}; this changes the operation
count by a constant factor and recovers the \(F\)-valued outputs.
The number field, basis and constants may depend on \(\varepsilon\)
and \(F\), but not on \(n\).  Thus \(w_F(k)\le2\); the output-size
lower bound gives equality, and \(\alpha_F\ge k>0.465\).
\end{proof}

The accuracy-dependent schemes in this corollary do not supply a common
finite set of exceptional positive characteristics.  No transfer of
\(\alpha_F>0.465\) to positive characteristic is asserted.
